% IV. Програмна реализация.
% Обща структура на изходния текст, йерархия на типовете, коментар на ключовите места.
% Листингите се вмъкват с \lstinputlisting от src/, за да не се разминават с кода.
\section{Програмна реализация}
\label{sec:implementation}

Изходният текст е на C++20 и се изгражда с CMake. Разделен е на заглавни файлове с
публичните типове и на единици за превод с реализацията, така че смяна на реализацията
на индекса да не води до преизграждане на целия проект. Векторизираната част е събрана
в отделни файлове, за да може да бъде компилирана с различни настройки от останалия
код и да бъде изключена изцяло при сравнение със скаларната реализация.

Ядрото на реализацията е заявката за най-близки съседи. При k-d дървото обхождането
поддържа купчина с текущите най-добри кандидати и отсича поддървета, чието разстояние
по разделящата равнина вече надхвърля най-лошото прието разстояние. Векторизацията се
прилага на нивото на листата: точките в един лист се проверяват на групи, вместо
поотделно, което е възможно именно заради разположението по компонента, описано в
Раздел~\ref{sec:design}. Размерът на листа става настройка, защото балансира дълбочината
на обхождане срещу дължината на векторизираната част.

Стъпките на обработката ползват индекса през общия му интерфейс. Филтрирането по
статистика на разстоянията до съседите, оценяването на нормали чрез анализ на главните
компоненти на околността, извличането на описания на околност~\cite{rusu2009fpfh},
сегментацията чрез нарастване на области по съгласуваност на нормалите и подгонването
на равнина с устойчиво оценяване~\cite{fischler1981ransac} са реализирани като отделни
стъпки с еднакъв договор. Съвместяването следва схемата на итеративната най-близка
точка~\cite{besl1992icp} с метрика до точката и с отхвърляне на съответствията над
зададен праг~\cite{rusinkiewicz2001icp}.

\TODO{йерархия на типовете (Фиг.), с легенда}

\TODO{листинг на заявката за $k$ най-близки съседи, вмъкнат с \textbackslash lstinputlisting
от src/, след като реализацията бъде завършена}

\TODO{описание на модулите за вход и изход и на поддържаните формати на файловете}
